\section{The result and the denominator improvement}\label{sec:introduction}
For $s>1$, let
\[
 \beta(s)=\sum_{t=0}^{\infty}\frac{(-1)^t}{(2t+1)^s}
\]
be the Dirichlet beta function. Its values at positive odd integers are
nonzero rational multiples of powers of $\pi$. The arithmetic of its
even values is much less understood. We prove the following finite-set
result.

\begin{theorem}\label{thm:main}
At least one of $\beta(4),\beta(6),\ldots,\beta(32)$ is irrational.
\end{theorem}

Krisman \cite[Theorem 1.1]{krisman} proves the corresponding assertion
through $\beta(36)$, with seventeen candidate values. We reduce the
number of candidates to fifteen by improving the denominator estimate
in his construction and then choosing new parameters.

We construct a sequence of linear forms in
$1,\beta(4),\beta(6),\ldots,\beta(32)$ with rational coefficients.
If all fifteen beta values were rational, one fixed integer would
clear their denominators. Multiplying the forms by that integer and
by factors clearing their coefficients would produce a sequence of
integers. We show that this sequence tends to zero but has infinitely
many nonzero terms, which is impossible.

\subsection{Where the saving occurs}
Krisman's rational functions have integer zeros and half-integer poles.
Their second derivatives produce sequences of linear forms in even
beta values beginning at $\beta(4)$. The denominator-clearing factors
must be small enough for the cleared forms still to tend to zero.

We call the integer $a$ the address of the pole $-a-\tfrac12$.
At an odd prime $p$, a local denominator comes from differences between
addresses that agree modulo $p$. If an address is alone in its residue
class, none of its principal-part coefficients has a denominator divisible
by $p$.
The rational constant of the linear form also contains evaluation
denominators. At a fixed evaluation point, however, only one residue
class can have a denominator divisible by $p$. A zero identity expresses
the same contribution using poles on either side of that point. When
the troublesome class contains just one pole, we choose the side that
avoids it.

Thus only classes containing at least two pole addresses enter the
large-prime cost. Krisman's argument already removes primes larger
than the pole span \cite[Section 4.3]{krisman}; retaining this information
also for smaller primes provides the saving used here. A related
exceptional-group cancellation appears in
\cite[Section 3.2]{finite}. The refinement here is the use of isolated
pole classes in the explicit denominator estimate.

\subsection{Construction, nonvanishing, and sources}
The rational kernel and its local valuation framework follow
\cite[Sections 2--6]{krisman}. The second-derivative device adapts
Rivoal's construction for odd zeta values \cite[Section 2]{rivoal}.
The role of zeros throughout the pole region in controlling rational
constants is also present in Zudilin's work
\cite[proof of Lemma 1]{zudilin}.

Smallness of the forms does not ensure that they are nonzero.
We use the prime-dilation argument of \cite[Section 4]{powered},
which implements the auxiliary-prime principle of Lai and Sprang
\cite[Lemma 2.3]{laisprang}. At indices equal to a fixed base index
times a large prime, a rescaling makes every beta coefficient vanish
modulo that prime while the rational constant reduces to a fixed,
nonzero seed constant. Section~\ref{sec:nonvanishing} gives the
calculation for the present parameters.

We use $v_p$ for the valuation on $\Q$ normalized by $v_p(p)=1$,
with $v_p(0)=+\infty$. A rational number is \emph{integral at $p$}
if its denominator is prime to $p$, and is a \emph{unit at $p$} if
its numerator is also prime to $p$. Congruences are taken in the ring
$\Z_{(p)}$ of rationals integral at $p$. The proof otherwise uses
partial fractions, Stirling's formula, and the prime number theorem.
